% ============================================================
%  part3/ch18-defects.tex  —  Chapter 18: Defects
% ============================================================

\chapter{Defects, Walls, Strings, and Nodes}
\label{ch:defects}

\section{Why topology obstructs coarsening}

Coarsening removes interfaces.
But topology may prevent complete elimination.

\begin{definition}[Vacuum manifold]
  Let $\Phi : X \to M$ be an order parameter field.
  The \emph{vacuum manifold} is the target space $M$.
\end{definition}

The topology of $M$ determines what structures may persist.

\begin{definition}[Fundamental group]
  $\pi_1(M)$ classifies loops in $M$ up to continuous
  deformation (homotopy).
\end{definition}

\begin{definition}[Winding number]
  For $M = S^1$, the winding number of a loop $\gamma$ is
  \[
    n = \frac{1}{2\pi}\oint_\gamma d\theta \in \mathbb{Z}.
  \]
\end{definition}

\begin{proposition}[Winding number is conserved]
  The winding number is invariant under continuous deformation.
\end{proposition}

\begin{proof}
  Under deformation $\theta(s,t)$,
  $\frac{dn}{dt} = \frac{1}{2\pi}\oint d(\partial_t\theta) = 0$
  since the integral of an exact form around a closed loop
  vanishes.
\end{proof}

\begin{theorem}[Topology forces singularities]
  \label{thm:top-singular}
  A nonzero winding number implies a singularity inside the loop.
\end{theorem}

\begin{proof}
  If no singularity exists, the field is continuous inside the
  loop and can be contracted to a point, implying $n = 0$.
  Contradiction.
\end{proof}

\begin{TechNote}[Topology vs dynamics]
  Dynamics determines how a system evolves.
  Topology determines what evolution is allowed.
  Free-energy descent cannot remove a defect protected by a
  nontrivial homotopy class.
  This distinction is essential for RSVP: cosmic filaments
  may persist not because they are energetically favoured,
  but because topology prevents their elimination.
\end{TechNote}

\section{Defect classification}

\begin{center}
\begin{tabular}{@{}lll@{}}
\toprule
\textbf{Homotopy group} & \textbf{Defect type} & \textbf{Cosmic analogue} \\
\midrule
$\pi_0(M)$ & Domain walls & Void boundaries \\
$\pi_1(M)$ & Strings / vortices & Filaments \\
$\pi_2(M)$ & Monopoles & Cluster nodes \\
$\pi_3(M)$ & Textures & Large-scale flows \\
\bottomrule
\end{tabular}
\end{center}

\begin{proposition}[Defect energy]
  A vortex with winding number $n$ in two dimensions has energy
  \[
    E \sim 2\pi n^2 \log\frac{R}{a},
  \]
  where $R$ is the system size and $a$ is the core size.
\end{proposition}

\begin{proof}
  $|\nabla\Phi| \sim n/r$ for a vortex.
  Energy density $e \sim n^2/r^2$.
  Integrate: $E = \int_a^R (n^2/r^2)(2\pi r)\,dr
  = 2\pi n^2 \log(R/a)$.
\end{proof}

The logarithmic scaling explains why vortices are long-lived:
their energy grows without bound as $R \to \infty$.

\begin{theorem}[Total topological charge is conserved]
  Annihilation requires $n_1 + n_2 = 0$.
  Only opposite-sign defects may disappear.
  Total topological charge is conserved.
\end{theorem}

\begin{AdmPreview}
  The connection between topological defects and admissibility
  will be made precise in Chapter~\ref{ch:sheaf}.
  A nonzero cohomology class $H^1(X, \mathcal{A}) \neq 0$
  is the sheaf-theoretic analogue of a nontrivial homotopy
  class: it measures the obstruction to globally assembling
  admissible local data.
  Cosmic filaments as topological defects become cosmological
  horizons as sheaf obstructions.
\end{AdmPreview}

\begin{theorem}[Universality of coarsening structures]
  \label{thm:universality-coarsening}
  Systems possessing (1) local interactions,
  (2) energy descent, and (3) topological constraints
  generically evolve toward sparse networked structures
  characterised by defects, branching, and scale-dependent
  coherence.
\end{theorem}

\begin{proof}[Proof sketch]
  Chapters~12--18 provide examples: spin systems,
  phase-separating fluids, foams, river basins, ecological
  networks, synchronisation fields, cosmological structure.
  All satisfy $\frac{d\mathcal{F}}{dt} \leq 0$ with nontrivial
  geometric or topological constraints.
  The resulting structures exhibit common scaling laws and
  persistent defect networks.
\end{proof}

\begin{HolonomyNote}
  After traversing Part~III, the following invariant becomes
  visible: every chapter has been describing the same process.

  Ising domains, Allen--Cahn interfaces, Cahn--Hilliard
  coarsening, soap foams, river basins, neural synchronisation,
  and cosmic filaments are not analogies.
  They are instances of the same mathematical structure:
  free-energy descent with topological constraints.

  The holonomy of Part~III is:

  \principle{
    The universe is not building structure.\\[4pt]
    It is removing inadmissible structure.\\[4pt]
    What survives is what topology prevents from being removed.
  }

  This is the bridge into Part~IV, where the RSVP free-energy
  functional is introduced as the physical realisation of the
  same mathematics.
\end{HolonomyNote}

\WhatSurvived{Topological coarsening $\phi(0) \mapsto \phi(t\to\infty)$}{%
  Topological charge\\
  Defect network topology\\
  Universality class%
}{%
  Defect positions\\
  Interface geometries\\
  Sub-domain fluctuations%
}{%
  Total charge exactly.\\
  Defect network topology\\
  partially (annihilations).%
}{%
  $\ker$: all configurations\\
  with identical topological\\
  charge distribution.%
}
